\documentclass[12pt]{article}

\usepackage[letterpaper,left=1.5in,right=1in,top=1in,bottom=0.9in]{geometry}
\usepackage[T1]{fontenc}
\usepackage[utf8]{inputenc}
\usepackage{setspace}
\usepackage{microtype}
\usepackage{titlesec}

\setstretch{1.05}
\setlength{\parindent}{0pt}
\setlength{\parskip}{0.6em}

\titleformat{\section}{\normalfont\bfseries\Large}{\thesection}{0.6em}{}
\titleformat{\subsection}{\normalfont\bfseries\large}{\thesubsection}{0.6em}{}

\begin{document}

\begin{titlepage}
\begin{center}

\vspace*{2cm}

{\LARGE \textbf{FLOWER WARS}}\\[1.2cm]

{\large A Screenplay}\\[2.5cm]

{\normalsize
\textbf{Flyxion}\\
\vspace{0.6cm}
\today\\
\vspace{0.6cm}
Draft 01
}

\vfill

\end{center}
\end{titlepage}

\section*{Log line}

In the shadow of a rising empire and a collapsing alliance network, Tlacaelel---a noble son without a throne---re-engineers war into a rule-governed ritual of capture, turning annihilation into spectacle and forcing an entire civilization to confront the cost of keeping the sun, the state, and its enemies alive.

\section*{Opening Direction}

Darkness. Not silence, but distance. Drums are audible, faint and uneven, as though carried across water. The sound does not swell; it persists, patient and indifferent.

A thin band of light appears at the bottom of the frame. Dawn emerges slowly, revealing mist drifting across canals. The city of Tenochtitlan resolves not as spectacle but as geometry: causeways, stepped temples, water broken into deliberate lines.

The sun crests the horizon. No music announces it. The light simply arrives.

\section*{Scene I}

\textbf{EXT. CAUSEWAY OUTSIDE TENOCHTITLAN --- MORNING}

Warriors return along the causeway in small, irregular clusters. There is no triumph in their posture. Some carry damaged shields. Others walk empty-handed. A few guide captives by the arm, not bound, but clearly unable to leave.

The camera does not linger on wounds. It lingers on absence. There are fewer prisoners than expected.

Tlacaelel walks among them, unarmored. He is young, but not inexperienced. His attention is precise, almost clerical. He stops beside a warrior cleaning obsidian fragments from a macuahuitl.

\textbf{TLACAELEL}  
How many?

The warrior looks up, uncertain whether this is a test.

\textbf{WARRIOR}  
Few. They fought hard.

Tlacaelel nods, as though this confirms something already calculated.

They pass a body being prepared for removal. Tlacaelel does not look away, but neither does he study it. His gaze moves instead to the surrounding warriors, measuring their fatigue, their silence.

\textbf{TLACAELEL}  
Hard fighting leaves little behind.

The warrior does not respond. The drums continue in the distance, unchanged.

Tlacaelel steps aside as captives are led past him toward the city. One stumbles. Another steadies him without being asked. Tlacaelel watches this carefully.

\textbf{TLACAELEL}  
(to himself)  
A living enemy still speaks.

He turns back toward the city. The sun climbs higher. The drums do not quicken.

\section*{Scene II}

\textbf{INT. COUNCIL CHAMBER --- LATE MORNING}

The chamber is wide and low, its ceiling supported by carved wooden beams darkened with age and smoke. Light enters from a high opening, falling unevenly across woven mats and stone benches.

Elders sit in deliberate spacing. No one rushes to speak. Authority here is expressed through patience.

Tlacaelel stands near the edge of the chamber, not centered. He waits.

At the far end, the \textbf{EMPEROR} listens without interruption as a priest finishes speaking. The priest’s tone is careful, formal, and exhausted.

\textbf{PRIEST}  
The signs were sufficient. The offering was accepted. The sun rose.

A murmur of assent moves through the chamber, brief and restrained.

Tlacaelel steps forward only after the murmur fades.

\textbf{TLACAELEL}  
The sun rose because it always rises.

The chamber stiffens. This is not blasphemy, but it is close enough to require response.

\textbf{ELDER}  
Choose your words.

Tlacaelel inclines his head.

\textbf{TLACAELEL}  
The sun rose because the world was not emptied yesterday.

Silence returns, thicker than before.

\textbf{TLACAELEL}  
We fought well. Too well. Our enemies died where they stood. They will not return. Neither will their sons.

A different elder leans forward, intrigued despite himself.

\textbf{SECOND ELDER}  
You speak as though this is loss.

\textbf{TLACAELEL}  
It is waste.

The word lands heavily. Waste is not a moral accusation; it is an administrative one.

\textbf{PRIEST}  
The gods do not require efficiency.

\textbf{TLACAELEL}  
No. They require continuity.

He gestures, small and precise.

\textbf{TLACAELEL}  
A dead enemy feeds the sun once. A living enemy feeds it again and again. Through fear. Through obligation. Through return.

Murmurs again, this time uneasy.

\textbf{ELDER}  
You would spare them?

\textbf{TLACAELEL}  
I would keep them.

The Emperor finally speaks, voice level, neither approving nor condemning.

\textbf{EMPEROR}  
You propose war without endings.

Tlacaelel meets his gaze.

\textbf{TLACAELEL}  
I propose war with rules.

The chamber absorbs this. No one agrees. No one dismisses it.

Light shifts across the floor as the sun climbs. The drums outside continue, unchanged.

\section*{Scene III}

\textbf{INT. TRAINING YARD WITHIN THE TEMPLE PRECINCT --- AFTERNOON}

The yard is enclosed by stone walls that mute the city beyond. Sunlight falls cleanly across packed earth. This is not a place of spectacle but of repetition.

Young warriors train in pairs. Their movements are deliberate, almost restrained. Blows are pulled at the last instant. Balance matters more than force.

Tlacaelel observes from the shade. Beside him stands a \textbf{VETERAN CAPTAIN}, scarred, alert, skeptical.

A trainee executes a clean disarm. His opponent freezes, blade at his throat. There is a pause, waiting for instruction.

\textbf{CAPTAIN}  
Finish it.

The trainee hesitates. Tlacaelel raises a hand. The trainee steps back. Both men lower their weapons.

The captain exhales sharply.

\textbf{CAPTAIN}  
You are teaching them to stop.

\textbf{TLACAELEL}  
I am teaching them to win twice.

They walk along the edge of the yard.

\textbf{CAPTAIN}  
A fight ends when one man cannot stand.

\textbf{TLACAELEL}  
A fight ends when both men remember it.

The captain watches another pair grapple. One forces the other to the ground and holds him there, immobile, breathing hard. No one cheers.

\textbf{CAPTAIN}  
And when the enemy does not agree to remember?

Tlacaelel does not answer immediately.

\textbf{TLACAELEL}  
Then we teach them the rules by surviving them.

The captain studies him now, not as a noble or advisor, but as a strategist.

\textbf{CAPTAIN}  
You would announce this to our enemies?

\textbf{TLACAELEL}  
Yes.

\textbf{CAPTAIN}  
You would tell them when and where to meet us.

\textbf{TLACAELEL}  
Yes.

The captain shakes his head.

\textbf{CAPTAIN}  
That is not war.

Tlacaelel watches the trainees separate, nod to one another, and reset.

\textbf{TLACAELEL}  
It is something harder. It requires restraint.

He turns back toward the yard.

\textbf{TLACAELEL}  
If we kill them all, they vanish.  
If we fight them this way, they return.  
If they return, they acknowledge us.

The captain remains silent. The training continues, rhythmic and controlled, like breath.

\section*{Scene IV}

\textbf{EXT. NEUTRAL FIELD BETWEEN CITY-STATES --- LATE AFTERNOON}

The field lies flat and open, deliberately unremarkable. Low grass, packed earth, no walls, no markers of ownership. Distance itself is the boundary.

Delegations from opposing city-states arrive from different directions, each stopping well short of the center. No one crosses the empty space between them.

Tlacaelel stands with a small escort. He carries no weapon. Across the field, an \textbf{ENVOY} from a rival city confers quietly with his attendants, then steps forward alone.

The space between them is measured by walking, not rushed.

They stop several paces apart.

\textbf{ENVOY}  
You ask us to meet without shields.

\textbf{TLACAELEL}  
I ask you to meet without surprise.

The envoy studies him carefully.

\textbf{ENVOY}  
Your messengers say you want war.

\textbf{TLACAELEL}  
Yes.

A pause.

\textbf{ENVOY}  
But not land.

\textbf{TLACAELEL}  
Not land.

\textbf{ENVOY}  
Not tribute.

\textbf{TLACAELEL}  
Not grain or gold.

The envoy allows himself a thin, incredulous smile.

\textbf{ENVOY}  
Then explain why we should listen.

Tlacaelel gestures toward the empty field.

\textbf{TLACAELEL}  
Because this is cheaper than annihilation.

The word hangs in the air, plain and unadorned.

\textbf{TLACAELEL}  
We will meet here.  
At times agreed upon.  
With warriors named in advance.  

He speaks as though reciting terms already written.

\textbf{TLACAELEL}  
Victory will be counted in captives, not bodies.  
Those taken will live.  
Those who fall will be few.

The envoy’s expression hardens.

\textbf{ENVOY}  
And if we refuse?

\textbf{TLACAELEL}  
Then we return to the old way.  
And one of us disappears.

A long silence. Wind moves through the grass.

\textbf{ENVOY}  
You would turn war into ceremony.

\textbf{TLACAELEL}  
I would turn it into memory.

The envoy looks back at his people, then again at the open field, imagining it filled.

\textbf{ENVOY}  
If we agree, this never ends.

Tlacaelel nods.

\textbf{TLACAELEL}  
That is the point.

The envoy considers this, then gives a single, deliberate nod.

\textbf{ENVOY}  
We will send names.

They step back from one another, retreating to their respective lines. The field remains empty, but no longer neutral. It has been designated.

The sun lowers toward the horizon. Shadows lengthen, orderly and calm.

\section*{Scene V}

\textbf{EXT. DESIGNATED FIELD --- MORNING}

The field has changed without being altered. Stones have been cleared. The grass is pressed flat in places by careful feet. Lines are implied by arrangement rather than drawn.

Warriors assemble on opposite sides, counted and named aloud by attendants. Each name is answered once. No chants. No music.

Tlacaelel stands apart from both groups, observing the symmetry. The warriors wear full regalia, but their weapons are newly prepared, edges honed with restraint rather than excess.

A priest approaches Tlacaelel quietly.

\textbf{PRIEST}  
The sun is high enough.

Tlacaelel nods.

\textbf{TLACAELEL}  
Begin.

There is no charge.

Pairs advance deliberately toward one another, seeking opponents rather than rushing the mass. Movements are precise, almost conversational. Each engagement isolates itself, briefly, before dissolving and reforming elsewhere.

The camera avoids impact. It follows hands gripping arms, feet hooking ankles, bodies unbalanced and guided to the ground.

A warrior from Tlacaelel’s side forces his opponent to kneel, blade poised, waiting. The opponent meets his gaze. Breath is visible in both chests.

The warrior does not strike. He signals. Attendants move in, separating them. The captured man is led away, alive.

This pattern repeats. Capture accumulates slowly, like a tally.

Tlacaelel watches, expression unreadable, but intent. The field hums with exertion, not chaos.

At the far edge, a younger warrior breaks form, raising his weapon too high, momentum threatening excess. A captain steps in instantly, restraining him, whispering something urgent. The moment collapses back into order.

By midmorning, the field holds more attendants than fighters. The fighting slows, then stops.

A horn sounds once. Clear. Final.

Warriors step back. Captives are counted. The sides remain distinct, but neither retreats in disorder.

Tlacaelel walks onto the field now, alone. He looks at the marks left in the dirt: impressions of knees, palms, heels. Evidence of struggle without erasure.

\textbf{TLACAELEL}  
(to no one)  
This will be remembered.

The sun continues its arc. The field empties, but its designation remains.

\section*{Scene VI}

\textbf{INT. TEMPLE ARCHIVE --- EVENING}

Scribes sit cross-legged before fresh codices. Their brushes move steadily. Images form: paired figures, repeated motifs, captives walking upright.

Tlacaelel stands behind them, overseeing without interference.

An elder approaches, subdued.

\textbf{ELDER}  
The people are confused.

\textbf{TLACAELEL}  
They will learn the rhythm.

\textbf{ELDER}  
They ask who won.

Tlacaelel considers this.

\textbf{TLACAELEL}  
Tell them the sun did.

The elder studies the codex, then Tlacaelel.

\textbf{ELDER}  
And if one day the sun asks for more?

Tlacaelel turns back to the images being recorded.

\textbf{TLACAELEL}  
Then we will already know how to give without ending ourselves.

The scribes continue. The codex grows. What was an event becomes doctrine.

\section*{Scene VII}

\textbf{INT. WARRIORS’ HALL --- NIGHT}

The hall is long and narrow, lit by evenly spaced braziers. Weapons rest against the walls, cleaned and ordered. Warriors sit in small clusters, speaking quietly. The atmosphere is unsettled, not angry.

A \textbf{YOUNG WARRIOR} sharpens his blade with unnecessary force. Across from him, an \textbf{OLDER WARRIOR} watches with mild disapproval.

\textbf{YOUNG WARRIOR}  
I had him. He was finished.

\textbf{OLDER WARRIOR}  
You had him alive.

The young warrior scoffs.

\textbf{YOUNG WARRIOR}  
That is not what I was trained for.

The older warrior gestures toward the hall.

\textbf{OLDER WARRIOR}  
You were trained to return.

The young warrior pauses, unsettled by the answer.

\textbf{YOUNG WARRIOR}  
They will say we are weaker.

\textbf{OLDER WARRIOR}  
They already say we are disciplined.

Silence settles between them.

At the far end of the hall, Tlacaelel enters quietly. Conversations diminish without announcement. He does not take the center of the room. He walks the perimeter, listening.

The young warrior rises.

\textbf{YOUNG WARRIOR}  
We are told to hold back. To wait. To capture.

Tlacaelel stops, but does not face him yet.

\textbf{TLACAELEL}  
You are told to choose.

The young warrior hesitates.

\textbf{YOUNG WARRIOR}  
Choose what?

Tlacaelel turns now.

\textbf{TLACAELEL}  
Whether you want one moment of certainty or many moments of consequence.

The hall is quiet enough to hear the fire shift.

\textbf{TLACAELEL}  
Killing ends a question.  
Capture keeps it open.

He moves again, leaving the young warrior standing, uncertain but attentive.

\section*{Scene VIII}

\textbf{INT. PRIESTS’ ENCLOSURE --- SAME NIGHT}

The enclosure is circular, enclosed, heavy with incense. Priests sit facing inward, speaking in measured tones.

A \textbf{SENIOR PRIEST} addresses the group.

\textbf{SENIOR PRIEST}  
The rites depend on certainty.

\textbf{SECOND PRIEST}  
The captives arrive alive. The sequence remains intact.

\textbf{SENIOR PRIEST}  
For now.

Tlacaelel stands at the threshold, visible but not entering.

\textbf{SENIOR PRIEST}  
You have changed the order of approach.

\textbf{TLACAELEL}  
I have changed the rate.

The priests exchange looks.

\textbf{SENIOR PRIEST}  
You risk teaching the people that the gods can be managed.

Tlacaelel considers this carefully.

\textbf{TLACAELEL}  
We have always managed them.  
We just pretended otherwise.

A murmur ripples, half alarmed, half relieved.

\textbf{SECOND PRIEST}  
If this fails—

\textbf{TLACAELEL}  
—then we will fail slowly enough to adapt.

The senior priest studies him, recognizing the deeper threat.

\textbf{SENIOR PRIEST}  
You are making the gods patient.

\textbf{TLACAELEL}  
I am making the people survivable.

Tlacaelel steps back, leaving the priests to their unease.

\section*{Scene IX}

\textbf{EXT. CITY OVERLOOKING THE CANALS --- DAWN}

The city wakes again. Canoes move through water. Markets assemble. Life proceeds without spectacle.

Tlacaelel stands on a terrace above the canals, watching ordinary motion reclaim the space once reserved for crisis.

Below him, two boys mimic the ritual contest in play, grappling, releasing, resetting. Neither seeks to end the game.

Tlacaelel observes this without satisfaction, only recognition.

\textbf{TLACAELEL}  
(softly)  
It has begun.

The sun rises. No blood marks the moment.

\section*{Scene X}

\textbf{EXT. TRAINING GROUNDS BEYOND THE CITY --- LATE MORNING}

The ground here is less formal than the temple yard. Grass breaks through packed soil. Wooden posts mark distances rather than boundaries. This is where new forms are tested before they are named.

A small troupe of \textbf{YOUNG WARRIORS} trains together, mixed in age and gender. Their armor is lighter, their movements faster, less ceremonial but no less disciplined. They move in coordinated patterns, rotating roles fluidly: attacker becomes defender, defender becomes observer.

A \textbf{YOUNG WOMAN} pivots low, sweeping an opponent’s leg and pinning him briefly before releasing. She does not celebrate. She resets.

Another woman signals, and the group reconfigures, circling inward, then outward, like breath.

Tlacaelel watches from a distance with the \textbf{VETERAN CAPTAIN}. Neither interrupts.

\textbf{CAPTAIN}  
They move differently.

\textbf{TLACAELEL}  
They learned after the rule was announced.

The captain studies the group more carefully now.

\textbf{CAPTAIN}  
They are not seeking the ending.

\textbf{TLACAELEL}  
They are seeking continuity.

A young man hesitates, uncertain, then yields space to a smaller opponent rather than overpowering her. She uses the opening to unbalance him cleanly. He laughs once, short and surprised, then steadies himself.

The captain frowns, then softens.

\textbf{CAPTAIN}  
This would not have happened before.

\textbf{TLACAELEL}  
Before, strength decided everything.  
Now, judgment does.

The troupe pauses as a horn sounds in the distance. They gather briefly, exchanging quiet words, adjusting straps, correcting one another’s stance without hierarchy.

One of the women looks toward Tlacaelel, meeting his gaze without deference or challenge, simply acknowledgment.

Tlacaelel inclines his head slightly.

\section*{Scene XI}

\textbf{INT. ARMORY ANNEX --- AFTERNOON}

Racks of weapons line the walls. New designs sit among older ones: modified grips, blunted edges near the spine, bindings meant to trap rather than cut.

The young troupe enters together. An \textbf{ARMORER} gestures as he explains changes.

\textbf{ARMORER}  
These are not weaker blades.  
They are more selective ones.

A young warrior tests the balance, nodding.

\textbf{YOUNG WOMAN}  
They require closer work.

\textbf{ARMORER}  
They require attention.

Tlacaelel stands near the doorway, listening.

\textbf{TLACAELEL}  
Attention is what keeps rules alive.

The group grows quiet, listening now not as trainees but as participants in doctrine.

\textbf{TLACAELEL}  
If this becomes spectacle alone, it will fail.  
If it becomes habit, it will survive.

He looks from face to face, measuring not loyalty but comprehension.

\textbf{TLACAELEL}  
You are not replacements for the old warriors.  
You are evidence that the war has changed.

The troupe absorbs this, not solemnly, but seriously.

\section*{Scene XII}

\textbf{EXT. SAME TRAINING GROUNDS --- SUNSET}

The light softens. The troupe trains once more, slower now, emphasizing balance and recovery. Movements echo those of the ritual field but are looser, adaptive.

In the distance, the city continues its rhythms uninterrupted.

Tlacaelel watches until the sun lowers enough that the forms blur together.

\textbf{TLACAELEL}  
(to himself)  
If they learn this first, they may never miss the other way.

The troupe disperses, laughing quietly, carrying their weapons with care rather than reverence.

The field remains, marked not by blood, but by repetition.

\section*{Scene XIII}

\textbf{INT. COUNCIL CHAMBER --- STORM AFTERNOON}

Rain strikes stone outside, irregular and loud enough to intrude. The chamber is fuller than before. Attendance signals urgency.

An \textbf{EMISSARY} kneels at the center, travel-worn, breathing controlled but tight.

\textbf{EMISSARY}  
They refused the field.

A ripple moves through the council, sharper than earlier murmurs.

\textbf{ELDER}  
Refused to meet?

\textbf{EMISSARY}  
Refused the terms. They attacked a caravan. No signal. No names exchanged.

Silence, heavy now with implication rather than thought.

The Emperor looks to Tlacaelel.

\textbf{EMPEROR}  
This is what you warned us against.

Tlacaelel inclines his head.

\textbf{TLACAELEL}  
Yes.

\textbf{ELDER}  
Then your system fails.

Tlacaelel does not answer immediately.

\textbf{TLACAELEL}  
It has been tested.

Rain intensifies, echoing through the chamber’s high opening.

\textbf{TLACAELEL}  
A rule does not fail because someone breaks it.  
It fails only if no one responds to the break.

\section*{Scene XIV}

\textbf{EXT. OUTER CAUSEWAY --- DUSK}

The causeway is tense, occupied by guards standing closer together than usual. Smoke rises faintly from beyond the lake, distant and unresolved.

Tlacaelel walks with the Veteran Captain.

\textbf{CAPTAIN}  
They will say this proves restraint invites chaos.

\textbf{TLACAELEL}  
Restraint reveals it.

They stop, watching messengers depart in different directions.

\textbf{CAPTAIN}  
What do you do when someone refuses the game?

Tlacaelel considers the water, dark and restless.

\textbf{TLACAELEL}  
You isolate them from the audience.

The captain studies him.

\textbf{CAPTAIN}  
You make an example.

\textbf{TLACAELEL}  
No.  
You make them irrelevant.

\section*{Scene XV}

\textbf{EXT. DESIGNATED FIELD --- DAYS LATER --- MORNING}

The field is prepared again, but something has changed. Attendance is larger. Observers gather at a distance, not celebrating, not chanting.

The troupe of young warriors stands among the assembled fighters, integrated rather than separate. Their posture is calm, alert.

Tlacaelel addresses the captains quietly.

\textbf{TLACAELEL}  
This field remains open only to those who accept return.

A horn sounds.

Opposing forces enter as before. Names are called. The ritual resumes.

At the far edge of the field, a \textbf{SCOUT} arrives at speed, breathless, whispering urgently to a captain. The captain’s eyes flick toward Tlacaelel.

Tlacaelel raises a hand.

The ritual does not stop.

Pairs engage. Capture proceeds. The audience watches carefully now, learning what endurance looks like.

Beyond the field, smoke still rises, distant but ignored.

\section*{Scene XVI}

\textbf{INT. TEMPLE TERRACE --- NIGHT}

Tlacaelel stands alone, overlooking the city. Torches mark steady points in the darkness.

The young woman from the troupe approaches, hesitating only briefly.

\textbf{YOUNG WOMAN}  
If they keep refusing—

Tlacaelel does not turn.

\textbf{TLACAELEL}  
Then they become noise.

She absorbs this.

\textbf{YOUNG WOMAN}  
And if the noise grows?

Tlacaelel finally faces her.

\textbf{TLACAELEL}  
Then we will discover whether the rule is strong enough to survive fear.

She nods once, not reassured, but resolved.

Below them, the city holds its shape.

\section*{Scene XVII}

\textbf{EXT. CITY AT DAWN --- TIME PASSING}

Morning returns. Markets open. Children rehearse ritual movements in play. Warriors train with measured force.

The Flower Wars continue, not uninterrupted, but persistent.

Tlacaelel watches the sun rise again, expression unchanged.

\textbf{TLACAELEL}  
(to himself)  
A system is not proven when it works.  
It is proven when it endures being refused.

The light settles. The city continues.

\section*{Scene XVIII}

\textbf{INT. TEMPLE ARCHIVE --- YEARS LATER --- DAY}

The archive has expanded. Codices line the walls in careful order. What were once single volumes are now families of record, variations layered over time.

Scribes work quietly. Their hands move with confidence born of repetition. Images of ritual combat recur with subtle differences: changing attire, shifting formations, new emblems marking alliances that have endured through cycles of contest.

Tlacaelel enters slowly. Age has marked him, but not diminished his precision. A \textbf{SENIOR SCRIBE} rises.

\textbf{SENIOR SCRIBE}  
The old battles are difficult to place now.

Tlacaelel studies a codex.

\textbf{TLACAELEL}  
That is because they ended.

He turns a page.

\textbf{TLACAELEL}  
These do not.

\section*{Scene XIX}

\textbf{EXT. RITUAL FIELD --- LATE AFTERNOON}

The field is familiar now. Permanent markers stand at its edges, not walls but acknowledgments. Seating has been arranged for observers, spaced to avoid spectacle.

The troupe of warriors has grown older. New members join seamlessly. The presence of women among them is unremarkable now, integrated into expectation rather than novelty.

Combat unfolds with practiced economy. Captures are efficient, almost inevitable.

An \textbf{ELDER OBSERVER} leans toward another.

\textbf{ELDER OBSERVER}  
We no longer count victories.

\textbf{SECOND OBSERVER}  
We count returns.

A captured warrior is released after formal acknowledgment. He walks back to his line upright, unshamed.

The field hums with recognition rather than tension.

\section*{Scene XX}

\textbf{INT. COUNCIL CHAMBER --- EVENING}

The council convenes in relative calm. Debate has softened into adjustment.

A \textbf{MERCHANT REPRESENTATIVE} speaks.

\textbf{MERCHANT}  
Trade routes have stabilized.

An elder nods.

\textbf{ELDER}  
The borders breathe.

Tlacaelel listens, hands folded.

\textbf{EMPEROR}  
You have given us war that behaves.

Tlacaelel responds evenly.

\textbf{TLACAELEL}  
I have given it a place to go.

\section*{Scene XXI}

\textbf{EXT. COASTLINE --- DAWN}

The horizon is unfamiliar. Large shapes interrupt the line between sea and sky. They move without rhythm, indifferent to tides.

Tlacaelel stands with scouts. The wind carries sounds no one recognizes.

\textbf{SCOUT}  
They do not send names.

Tlacaelel watches the distant vessels.

\textbf{TLACAELEL}  
They do not know the field.

\section*{Scene XXII}

\textbf{INT. COUNCIL CHAMBER --- URGENT}

Maps are spread across the floor. The markings do not align with known boundaries.

\textbf{ELDER}  
They take without contest.

\textbf{PRIEST}  
They do not wait for ritual.

Tlacaelel remains composed.

\textbf{TLACAELEL}  
Then they are not enemies in the old sense.

The Emperor studies him.

\textbf{EMPEROR}  
Can your system survive this?

Tlacaelel looks to the maps, to the unfamiliar symbols.

\textbf{TLACAELEL}  
It can survive many things.  
But it cannot teach those who refuse to see one another.

\section*{Scene XXIII}

\textbf{EXT. RITUAL FIELD --- EMPTY --- NIGHT}

The field lies unused. Moonlight outlines its markers. No footsteps disturb the ground.

Tlacaelel walks its perimeter alone, slower now. He pauses at the center.

\textbf{TLACAELEL}  
(quietly)  
A rule requires mutual recognition.

He looks outward, beyond the field, beyond the city.

\textbf{TLACAELEL}  
And recognition cannot be forced.

He turns back toward the city lights, which remain steady.

\section*{Scene XXIV}

\textbf{INT. TEMPLE TERRACE --- PRE-DAWN}

The troupe of warriors assembles, older leaders among younger ones. Women and men stand together without distinction.

Tlacaelel addresses them one final time.

\textbf{TLACAELEL}  
You were trained to return.  
Remember that, even if others do not.

They listen in silence, absorbing the weight of continuity.

The first light of dawn touches the city.

\section*{Scene XXV}

\textbf{EXT. CITY AT SUNRISE}

The sun rises once more. The city moves. The field remains.

The Flower Wars endure, not as guarantee, but as memory made structure.

Tlacaelel watches the light settle across water and stone.

\textbf{TLACAELEL}  
(to himself)  
We did not stop violence.  
We taught it where to live.

The sun climbs. The city breathes.

\section*{Epilogue}

\section*{Scene XXVI}

\textbf{INT. TEMPLE ARCHIVE --- NIGHT}

The archive is quiet in a way it has not been before. Fewer scribes remain. The codices are complete enough to stand without constant amendment.

An \textbf{APPRENTICE SCRIBE}, young and careful, traces a finished glyph with his finger. He hesitates, then looks up at Tlacaelel, now seated, older, his posture unchanged but slowed by time.

\textbf{APPRENTICE SCRIBE}  
Will this be remembered as war?

Tlacaelel considers the question without urgency.

\textbf{TLACAELEL}  
It will be remembered as structure.

\textbf{APPRENTICE SCRIBE}  
And if the structure breaks?

Tlacaelel gestures toward the shelves.

\textbf{TLACAELEL}  
Then what survives is not the rule, but the idea that rules were possible.

The apprentice nods, committing this to memory rather than ink.

\section*{Scene XXVII}

\textbf{EXT. RITUAL FIELD --- MANY YEARS LATER --- DAY}

The field is quieter now. Fewer contests are held. When they occur, they are smaller, more symbolic.

Children play at the edge, watched by elders who correct form gently, without urgency. A girl pins a boy and releases him immediately. He bows, exaggerated, and runs off laughing.

The markers remain upright.

An old woman, once among the early troupe of warriors, walks the perimeter with a staff. She stops, rests, and watches the children.

\textbf{OLD WOMAN}  
(to herself)  
They still know where it is.

She turns back toward the city.

\section*{Scene XXVIII}

\textbf{EXT. CAUSEWAY --- DUSK}

Merchants move slowly. Canoes pass without escort. Life is quieter, not safer, but less brittle.

A traveler pauses at the edge of the causeway, asking directions. No one reaches for a weapon.

\section*{Scene XXIX}

\textbf{INT. TEMPLE TERRACE --- NIGHT}

Tlacaelel stands once more at the terrace, supported now by a carved rail. The city lights below him are fewer but steady.

A priest approaches, respectful, restrained.

\textbf{PRIEST}  
They say the gods have changed.

Tlacaelel does not turn.

\textbf{TLACAELEL}  
No.  
They learned to wait.

The priest studies him, searching for regret.

\textbf{PRIEST}  
And you?

Tlacaelel looks out across water and stone.

\textbf{TLACAELEL}  
I learned that violence does not disappear.  
It learns.

The priest bows slightly and withdraws.

\section*{Scene XXX}

\textbf{EXT. RITUAL FIELD --- NIGHT}

The field lies empty under the stars. The markers cast long shadows.

Wind moves across the grass. Nothing else stirs.

Tlacaelel enters slowly, stands at the center, and closes his eyes.

\textbf{TLACAELEL}  
(softly)  
If the sun rises tomorrow, it will not be because of us alone.

He opens his eyes and looks upward.

\section*{Final Image}

The sun rises over the city one final time. Water reflects light. Stone holds warmth. The field remains visible, not central, not forgotten.

\textit{War did not end.  
It was taught where to stand.}

\section*{Historical Note}

\textit{Flower Wars} is inspired by historical traditions associated with the Mexica (Aztec) polity of the fifteenth century and by later scholarly reconstructions of the role played by Tlacaelel, a high-ranking statesman and military organizer whose influence on imperial ideology far exceeded his formal titles.

The \textit{xōchiyāōyōtl}, commonly translated as “Flower Wars,” are understood by historians as a form of prearranged conflict between allied or rival city-states, most notably involving Tenochtitlan, Tlaxcala, and Huexotzinco. These engagements were distinct from wars of conquest. They were governed by conventions regarding timing, location, and conduct, and emphasized the taking of captives rather than territorial annihilation. Captives served religious, political, and social functions, including ritual sacrifice, the reaffirmation of warrior status, and the maintenance of inter-polity relationships.

Tlacaelel is traditionally credited, in both Indigenous accounts and colonial-era chronicles, with reshaping Mexica cosmology and statecraft following a period of crisis in the early fifteenth century. Among these reforms were the elevation of Huitzilopochtli to a central position in state religion, the rewriting of dynastic histories, and the reorganization of warfare as a structural necessity rather than an episodic event. While modern scholars debate the extent to which these reforms can be attributed to a single individual, there is broad agreement that the Mexica state underwent a deliberate ideological consolidation during this period.

This screenplay adopts a restrained and interpretive approach to these traditions. Dialogue, rituals, and institutions depicted here are not verbatim reconstructions but dramatized extrapolations grounded in historical scholarship. The portrayal of women participating directly in warrior training reflects evidence that Mexica society recognized martial, symbolic, and ritual roles for women, even if formal battlefield participation remains debated. Their inclusion here emphasizes continuity, adaptability, and the social transmission of disciplined conflict rather than strict historical literalism.

Finally, the film resists framing the Flower Wars as either barbaric spectacle or moral allegory. Instead, it presents them as an early and sophisticated attempt to regulate violence within a political system that recognized the destructive limits of total war. The narrative neither endorses nor condemns this system, but treats it as a historical experiment in constraint: one whose logic functioned only so long as all participants recognized the field, the rules, and one another.

Any anachronisms are intentional, serving clarity rather than correction. The aim is not to modernize the past, but to render visible a historical moment in which war was deliberately redesigned as a repeatable, survivable institution.

\end{document}